% Compile with XeLaTeX. Requires Times New Roman and Latin Modern Math.
\documentclass[aspectratio=169,10pt]{beamer}
\usepackage{fontspec}
\usepackage{amsmath}
\usepackage{unicode-math}
\usefonttheme{serif}
\setmainfont{Times New Roman}
\setsansfont{Times New Roman}
\setmathfont{Latin Modern Math}
\setbeamertemplate{navigation symbols}{}
\setbeamertemplate{footline}{}
\setbeamertemplate{headline}{}
\setbeamertemplate{frametitle}{\vspace{3mm}\insertframetitle\par}
\setbeamercolor{normal text}{fg=black,bg=white}
\setbeamercolor{frametitle}{fg=black,bg=white}
\setbeamerfont{frametitle}{size=\large,series=\bfseries}
\setbeamersize{text margin left=8mm,text margin right=8mm}
\setlength{\parskip}{2pt}
\newcommand{\Db}{\overline D}
\newcommand{\Pb}{\overline P}
\newcommand{\rh}{\rho_0}
\begin{document}

\begin{frame}{1. Feifei's quorum sensing model}
Particles diffuse faster when their perceived crowding exceeds a threshold.
\[
 \mathrm d\mathbf r_i=\sqrt{2D(P_i)}\,\mathrm d\mathbf W_i\;,
 \qquad P_i=\sum_{\substack{j\ne i\\0<r_{ij}<d_c}}\frac{1}{2\pi r_{ij}}\;.
\]
\[
D(P)=\begin{cases}D_{\min},&P\le P_{\mathrm{th}},\\D_{\max},&P>P_{\mathrm{th}},\end{cases}
\qquad 0<D_{\min}<D_{\max}.
\]
\textbf{Nonlocal perception: an unnormalized, isotropic sensing kernel}
\[
K(r)=\frac{\Theta(d_c-r)}{2\pi r},\qquad P=K*\rho,
\qquad\int K(r)\,\mathrm d^2r=d_c.
\]
\textbf{Itô expansion and random-current covariance give the density equation}
\[
\partial_t\rho=\nabla^2[\rho D(P)]
 -\nabla\!\cdot\!\left[\sqrt{2\rho D(P)}\,\boldsymbol\Lambda\right].
\]
\small $\boldsymbol\Lambda$ is unit vector white noise. Self-sensing is excluded.
% Source: Section 3.1--3.2, Eqs. (74)--(79).
\end{frame}

\begin{frame}{2. Linear fluctuations and the structure factor}
\textbf{Assume a smooth response near the uniform state}
\[
\rho=\rh+\delta\rho,\qquad \Pb=\rh d_c,\qquad
\Db=D(\Pb)>0,\qquad g=D'(\Pb).
\]
\[
\partial_t\delta\rho=
\nabla^2\!\left[\Db\,\delta\rho+\rh g(K*\delta\rho)\right]
-\nabla\!\cdot\!\left[\sqrt{2\rh\Db}\,\boldsymbol\Lambda\right].
\]
\textbf{Each nonzero Fourier mode is a scalar Ornstein--Uhlenbeck process}
\[
\partial_t\delta\rho_{\mathbf k}=-\lambda(k)\delta\rho_{\mathbf k}+\eta_{\mathbf k},
\qquad\lambda(k)=k^2\left[\Db+\rh g\widehat K(k)\right].
\]
\[
\underbrace{Q(k)=2\rh\Db k^2}_{\text{noise power}},\qquad
\underbrace{2\lambda(k)C(k)=Q(k)}_{\text{stationary variance}},\qquad
C(k)=\langle|\delta\rho_{\mathbf k}|^2\rangle.
\]
\[
\mathord{\displaystyle S(k)=\frac{C(k)}{\rh}
=\frac{\Db}{\Db+\rh g\widehat K(k)}}
\]
\small Valid for stable modes, $\Db+\rh g\widehat K(k)>0$, within the linear-noise approximation.
% Source: Section 3.2, Eqs. (80)--(86). Volume-normalized Fourier convention.
\end{frame}

\begin{frame}{3. Why an approximately linear spectrum appears}
\textbf{The 2D Fourier transform reduces to a Bessel integral}
\[
\widehat K(k)=\int_0^{d_c}J_0(kr)\,\mathrm dr
=\frac{F(kd_c)}{k},\qquad F(x)=\int_0^x J_0(u)\,\mathrm du.
\]
Define $k_*={\rh g}/{\Db}>0$ and $B=k_*d_c$. Then
\[
\displaystyle S(k)=\left[1+\frac{k_*}{k}F(kd_c)\right]^{-1}.
\]
For $kd_c\gg1$, $F(kd_c)\simeq1$ and $\widehat K(k)\simeq1/k$.
\[
\mathord{\displaystyle \underbrace{B\gg1}_{\text{For $P$ \rightarrow $P_{\mathrm{th}}$, $D(P)$ has a sharp transition $g \rightarrow \infty$, $B \rightarrow \infty$}},\quad d_c^{-1}\ll k\ll k_*
\quad\Longrightarrow\quad S(k)\simeq\frac{k}{k_*}}
\]
\textbf{Mechanism: the sensing kernel changes the relaxation scaling}
\[
\lambda(k)\simeq\rh gk,\qquad Q(k)\propto k^2,
\qquad S(k)=\frac{Q(k)}{2\rh\lambda(k)}\propto k.
\]
\small This intermediate range must lie below the microscopic and continuum cutoffs.
% Source: Section 3.3, Eqs. (87)--(92).
\end{frame}

\begin{frame}{4. Finite-range limit and the threshold approximation}
\textbf{At fixed finite sensing radius, the true small-$k$ limit is a plateau}
\[
\widehat K(k)=d_c\!\left[1-\frac{(kd_c)^2}{12}+\cdots\right],
\qquad \lim_{k\to0}S(k)=\frac{1}{1+B}>0.
\]
A finite-range $S(k)\sim k$ window alone does not establish asymptotic hyperuniformity.
\medskip

\textbf{The microscopic step needs an explicit smooth or effective response}
\[
D_\varepsilon(P)=D_{\min}+
\frac{D_{\max}-D_{\min}}{1+\exp[-(P-P_{\mathrm{th}})/\varepsilon]}.
\]
\[
g=D_\varepsilon'(\Pb)>0\text{ near the switching region},
\qquad |\delta P|\ll\varepsilon.
\]
For the literal step, $g=0$ away from threshold and is undefined at threshold.
\medskip

Thus the prediction is a smooth-response or effective-closure result.
% Source: Section 3.3--3.4, Eqs. (93)--(99). Conclusions refer to homogeneous linear theory.
\end{frame}
\end{document}
